% ==============================================================================
% PREAMBLE COMMUN - Certification ECNSA
% ------------------------------------------------------------------------------
% Ce fichier ne contient pas de texte : il ne fait que charger des paquets et
% régler la mise en page. Il doit être appelé APRÈS \documentclass et AVANT
% \begin{document}.
%
% Rappel : un « paquet » (package) est un fichier de commandes LaTeX fourni par
% une distribution (TeX Live, MiKTeX...). \usepackage{nom} active ses commandes.
% ==============================================================================

% ------------------------------------------------------------------------------
% 1. ENCODAGE ET LANGUE
% ------------------------------------------------------------------------------

% inputenc : indique à LaTeX comment lire le fichier source.
% L'option [utf8] permet d'écrire directement les accents (é, à, ç...).
\usepackage[utf8]{inputenc}

% fontenc : choisit l'encodage des polices de sortie.
% L'option [T1] améliore le rendu des accents et le copier-coller depuis le PDF.
\usepackage[T1]{fontenc}

% babel : applique les règles typographiques d'une langue.
% L'option [french] gère les espaces avant : ; ! ?, les guillemets français, etc.
\usepackage[french]{babel}

% lmodern : charge la famille de polices Latin Modern.
% Elle est plus lisible et plus complète que la police par défaut.
\usepackage{lmodern}

% ------------------------------------------------------------------------------
% 2. MISE EN PAGE (GEOMETRY)
% ------------------------------------------------------------------------------

% geometry : règle les marges et le format du papier.
% Réglage par défaut ; chaque profil le redéfinit si nécessaire.
\usepackage[a4paper, margin=2cm]{geometry}

% ------------------------------------------------------------------------------
% 3. TABLEAUX
% ------------------------------------------------------------------------------

% booktabs : fournit les traits horizontaux professionnels des tableaux :
%   \toprule    -> trait épais du haut
%   \midrule    -> trait fin après l'en-tête
%   \bottomrule -> trait épais du bas
\usepackage{booktabs}

% array : améliore la définition des colonnes de tableau.
% Permet d'appliquer un format à une colonne :
%   >{\\bfseries}p{3cm}  -> colonne de 3 cm en gras
%   >{\\raggedright\\arraybackslash}p{3cm} -> colonne alignée à gauche
\usepackage{array}

% tabularx : crée des tableaux qui occupent toute la largeur imposée.
% Sa colonne spéciale X se partage automatiquement l'espace restant.
\usepackage{tabularx}

% Colonne Y = colonne X alignée à gauche (utilisée dans les tableaux originaux).
\newcolumntype{Y}{>{\raggedright\arraybackslash}X}

% ------------------------------------------------------------------------------
% 4. IMAGES ET GRAPHIQUES
% ------------------------------------------------------------------------------

% graphicx : charge les images avec \includegraphics.
\usepackage{graphicx}

% caption : personnalise les légendes des figures et des tableaux.
\usepackage{caption}

% float : fournit l'option [H] pour les environnements figure et table.
\usepackage{float}

% xcolor : permet d'utiliser des couleurs si ponctuellement nécessaire.
\usepackage{xcolor}

% ------------------------------------------------------------------------------
% 5. NUMÉROTATION DES PAGES
% ------------------------------------------------------------------------------

% pagestyle plain : numéro de page centré en bas, pas d'en-tête.
% C'est le réglage retenu pour tous les types de documents.
\pagestyle{plain}

% ------------------------------------------------------------------------------
% 6. LISTES ET ESPACEMENTS
% ------------------------------------------------------------------------------

% enumitem : règle finement les listes itemize et enumerate.
\usepackage{enumitem}

% titlesec : personnalise les espacements autour des titres de sections.
\usepackage{titlesec}

% \parindent : supprime l'alinéa de première ligne (style documents techniques).
\setlength{\parindent}{0pt}

% ------------------------------------------------------------------------------
% 7. LIENS ET NAVIGATION
% ------------------------------------------------------------------------------

% hyperref : rend cliquables les liens, le sommaire et les renvois.
% [hidelinks] supprime les cadres colorés autour des liens.
\usepackage[hidelinks]{hyperref}

% ------------------------------------------------------------------------------
% 8. RÉGLAGES GLOBAUX PAR DÉFAUT
% ------------------------------------------------------------------------------

% \arraystretch : hauteur des lignes de tableau.
% 1.25 ajoute 25 % de hauteur pour éviter que le texte touche les traits.
\renewcommand{\arraystretch}{1.25}

% \captionsetup : légendes en petite taille, label en gras.
\captionsetup{font=small, labelfont=bf}

% Évite les débordements dans les URL, les lignes de code et les tableaux.
\setlength{\emergencystretch}{2em}
\urlstyle{same}
